% R&R Other Project Information, Field 11: Equipment
% (NOFO N0001426SF004, pp. 20-21), for the copper/DED alternative
% (onr-yip-technical-proposal-cu-ded.tex). Same as field11-equipment.tex
% except the deposition process and the powder checks.
% Build: xelatex field11-equipment-cu-ded.tex
\documentclass[12pt,letterpaper]{article}
\usepackage[margin=1in]{geometry}
\usepackage{fontspec}
\IfFontExistsTF{Times New Roman}{\setmainfont{Times New Roman}}{\setmainfont{TeX Gyre Termes}}
\usepackage{microtype}
\usepackage{xcolor}
\usepackage{xurl}
\usepackage[colorlinks=true,urlcolor=blue!70!black]{hyperref}
\setlength{\parindent}{0pt}
\setlength{\parskip}{5pt}
\pagestyle{plain}
\newcommand{\todo}[1]{\textcolor{red}{[TODO Sterling: #1]}}

\begin{document}
\begin{center}\textbf{Equipment}\\
ONR FY2027 YIP, NOFO N0001426SF004. PI: Sterling G.\ Baird, Brigham Young University\end{center}

\textbf{Equipment available.} The equipment below is in place and used at no
cost to the project except where a facility charges hourly rates (shown in
the budget).
\begin{itemize}\setlength{\itemsep}{1pt}
\item AMAZEMET rePOWDER ultrasonic atomizer with induction crucible module
(PI laboratory, BYU; bought with start-up funds).
\item Liquid-handling gantry (CubXL, Ursa Laboratories) with a pipette, a
vial capper, and cameras, run by scripts (PI laboratory); it runs the
corrosion screen and samples immersion solutions.
\item Powder handling, sieving, and specimen preparation equipment in the PI
laboratory \todo{list the main items; add a Hall flowmeter funnel if one
must be bought}.
\item SEM, EDS, EBSD, and TEM at the BYU Microscopy Center (Apreo C,
Helios 600, Verios; Tecnai TF-20).
\item Agilent 8900 ICP-MS/MS and XPS (BYU Spectroscopic Facility); X-ray
diffraction (BYU XRD Facility).
\item Metallographic polishers, a Keyence VHX-7000 digital microscope, and
CNC machining (BYU).
\item Tensile testing: \todo{BYU load frame or test service}.
\item BYU high-performance computing cluster.
\item Blown-powder laser DED system at \todo{facility; candidate: University
of Utah, Prof.\ Zak Fang}, used through purchased sessions \todo{system model;
likely an InssTek MX-Lab}.
\end{itemize}

\textbf{Equipment to be bought with project funds.} One potentiostat with
corrosion software and a flat cell, Year 1, estimated at \$18,000
\todo{attach the vendor quote, as the NOFO requires}. The electrochemical
corrosion screen is one step of the workflow that the search method
schedules, so the instrument must be available every day and run by the
laboratory's own scripts on the gantry above. No equipment that can serve
this role is available to the project at BYU \todo{decide: a module
potentiostat with impedance on the gantry (about \$600 in the lab's survey)
plus a borrowed laboratory instrument for cross-checks would remove this
item and most of the budget line; see cubxl-electrochemistry.md}.

\textbf{Government-owned equipment.} None.

\end{document}
